\section{A finite-adelic image theorem}\label{sec:image}

For a prime $p$, write $\Z_p$ for the $p$-adic integers, $v_p$ for the
normalized valuation, and $|z|_p=p^{-v_p(z)}$ for $z\neq0$. Set
$v_p(0)=+\infty$. If $s,t\geq0$, define
\begin{equation}\label{eq:kernel}
 \Hp{p}{s}{t}(z)=p^{-s v_p(z)-t p^{v_p(z)}}\qquad(z\neq0),
\end{equation}
and prescribe
\begin{equation}\label{eq:kernel_zero}
 \Hp{p}{s}{t}(0)=
 \begin{cases}
  0,&(s,t)\neq(0,0),\\
  1,&(s,t)=(0,0).
 \end{cases}
\end{equation}
An exponent pair is called \emph{active} when $(s,t)\neq(0,0)$.
Every kernel is continuous and takes values in $[0,1]$.
Indeed, the valuation is locally constant away from zero, and an active
kernel tends to zero as the valuation tends to infinity.
The inactive kernel $\Hp{p}{0}{0}$ is the constant function one.
Convention~\eqref{eq:kernel_zero} avoids any interpretation of $0^0$.

Fix the following data throughout Sections~\ref{sec:image}--\ref{sec:perfectness}:
\begin{itemize}
 \item a finite set $S$ of primes and a real number $\Lambda>1$;
 \item a positive periodic sequence $(r_n)_{n\in\Z}$;
 \item for every $p\in S$, nonnegative real periodic sequences
 $(s_{p,n})_{n\in\Z}$ and $(t_{p,n})_{n\in\Z}$ with a common period
 $m_p$ satisfying $\gcd(m_p,p)=1$.
\end{itemize}
Let $w\geq1$ be a common period of all the sequences. A prime is active if
at least one of its exponent sequences is not identically zero. Inactive
primes can be removed from $S$ without changing any of the functions below.
We make this removal and write $d=|S|$. The list of exponent pairs at
each prime is finite, but a pair is not required to have integer or rational
entries.

Give $\Z/w\Z$ the discrete topology and set
\begin{equation}\label{eq:detector_definition}
 \D=\overline{\{(n\bmod w,(n)_{p\in S}):n\in\Z\}}
 \ \subseteq\ (\Z/w\Z)\times\prod_{p\in S}\Z_p.
\end{equation}
This is a nonempty compact metrizable group. On it define
\begin{equation}\label{eq:detector_function}
 \Phi(a,x)=\sum_{\ell=0}^{\infty}\Lambda^{-\ell}r_{a-\ell}
       \prod_{p\in S}\Hp{p}{s_{p,a-\ell}}{t_{p,a-\ell}}(x_p-\ell).
\end{equation}
The indices $a-\ell$ are evaluated by periodicity, so the summand does not
depend on a representative of $a\in\Z/w\Z$. Write
\begin{equation}\label{eq:RB}
 R=\max_n r_n,\qquad B=\frac{R}{1-\Lambda^{-1}}.
\end{equation}
Each summand of~\eqref{eq:detector_function} is bounded by
$R\Lambda^{-\ell}$. The Weierstrass criterion therefore gives uniform
convergence and continuity of $\Phi$.

For a nonempty compact $E\subset\R$, let $N_{\eps}(E)$ denote the least
number of closed real intervals of length $\eps$ covering $E$. The upper
box dimension is
\begin{equation}\label{eq:box_definition}
 \boxdim E=\limsup_{\eps\downarrow0}
          \frac{\log N_{\eps}(E)}{\log(1/\eps)}.
\end{equation}
The term \emph{Cantor set} will mean a nonempty compact, perfect,
totally disconnected metric space, here realized as a nowhere dense subset
of $\R$. Such a space is homeomorphic to the usual Cantor space.

\begin{theorem}[Quantitative image theorem]\label{thm:image}
For the fixed data above, the following statements hold.
\begin{enumerate}
 \item If $d=0$, then $\Phi(\D)$ is the finite set
 \begin{equation}\label{eq:finite_image}
  \left\{\frac{\displaystyle\sum_{\ell=0}^{w-1}
              r_{a-\ell}\Lambda^{-\ell}}
             {1-\Lambda^{-w}}:a\in\Z/w\Z\right\}.
 \end{equation}
 \item If $d>0$, then $\Phi(\D)$ is a Cantor set.
 \item If $d>0$, there are $C,\eps_0>0$, depending on the fixed data,
 such that
 \begin{equation}\label{eq:main_cover}
  N_{\eps}\bigl(\Phi(\D)\bigr)
   \leq C\bigl(1+\log(1/\eps)\bigr)^{2d}
   \qquad(0<\eps<\eps_0).
 \end{equation}
 Consequently $\boxdim\Phi(\D)=\hausdim\Phi(\D)=0$.
\end{enumerate}
\end{theorem}

Only positivity and periodicity are assumed for $r_n$, not the stronger
gcd-sequence property used in the dynamical application. The coprimality
of the exponent periods with their own prime is used to prove perfectness.
The covering estimate does not need that coprimality. These roles will be
kept separate in the proof.
